\documentclass[a4paper,11pt]{article}
\usepackage[a4paper,margin=2cm]{geometry}
\usepackage[T1]{fontenc}
\usepackage[utf8]{inputenc}
\usepackage[ngerman,english]{babel}
\usepackage{lmodern,microtype,xcolor,parskip,enumitem}
\usepackage[most]{tcolorbox}
\usepackage{titlesec}
\usepackage[hidelinks]{hyperref}
\definecolor{HeaderBlueDark}{RGB}{70,110,170}
\definecolor{RowBlueLight}{RGB}{235,243,255}
\definecolor{RowBlueDark}{RGB}{220,233,250}
\definecolor{SoftGray}{RGB}{90,90,90}
\setlength{\parindent}{0pt}
\setlist[itemize]{leftmargin=1.4em,itemsep=0.35em,topsep=0.35em}
\titleformat{\section}[block]{\Large\bfseries\color{HeaderBlueDark}}{}{0pt}{}
\newtcolorbox{exbox}{enhanced,breakable,colback=RowBlueLight,colframe=RowBlueDark,boxrule=0.4pt,arc=1.5mm,left=1.2mm,right=1.2mm,top=1mm,bottom=1mm,title={Beispiele \textnormal{(Examples)}},fonttitle=\bfseries,colbacktitle=RowBlueDark,coltitle=black}
\NewDocumentEnvironment{grammarentry}{mm+b}{\begin{tcolorbox}[enhanced,breakable,colback=white,colframe=HeaderBlueDark,title={#1\hfill\normalfont\footnotesize #2},colbacktitle=HeaderBlueDark,coltitle=white,fonttitle=\bfseries]#3\end{tcolorbox}}{}
\newcommand{\GermanText}[1]{\textbf{Deutsch: }#1\par}
\newcommand{\EnglishText}[1]{{\small\color{SoftGray}\textbf{English: }#1\par}}
\newcommand{\ExampleItem}[2]{\item \textbf{#1}\\{\small\color{SoftGray}#2}}
\title{Die Modalpartikel \glqq vielleicht\grqq\\{\large\textnormal{The modal particle vielleicht}}}
\date{}
\begin{document}
\maketitle
\tableofcontents
\section{Bedeutung und Gesprächsfunktion \textnormal{(Meaning and pragmatic function)}}
\begin{grammarentry}{vielleicht}{Modalpartikel / Abtönungspartikel}
\GermanText{Vielleicht kann in Ausrufen eine starke emotionale Bewertung oder überraschtes Erstaunen ausdrücken.}
\EnglishText{Vielleicht can convey a strong emotional evaluation or surprise in exclamations.}
\end{grammarentry}
\section{Satzart, Stellung und Betonung \textnormal{(Sentence mood, position and stress)}}
\begin{grammarentry}{Verwendung}{Modalpartikel / Abtönungspartikel}
\GermanText{Als Abtönungspartikel besonders in emotionalen Ausrufen und nicht im gewöhnlichen Sinne von möglicherweise.}
\EnglishText{As a modal particle especially in emotional exclamations rather than with the ordinary meaning perhaps.}
\end{grammarentry}
\GermanText{Als Abtönungspartikel ist vielleicht unflektierbar und verlangt keinen Kasus. Sie steht normalerweise im Mittelfeld und kann das Vorfeld eines normalen Aussagesatzes nicht allein besetzen.}
\EnglishText{As a discourse particle, vielleicht is invariable and governs no case. It normally occurs in the middle field and cannot stand alone in the prefield of a declarative sentence.}
\section{Typische Beispiele \textnormal{(Typical examples)}}
\begin{exbox}\begin{itemize}
\ExampleItem{Das war vielleicht ein Chaos!}{What a mess that was!}
\ExampleItem{Das war vielleicht eine Überraschung!}{What a surprise that was!}
\ExampleItem{Du bist vielleicht ein Scherzkeks!}{You're quite the joker! (possibly ironic)}
\ExampleItem{Das war vielleicht anstrengend!}{Boy, that was exhausting!}
\end{itemize}\end{exbox}
\section{Andere grammatische Rollen \textnormal{(Other grammatical roles)}}
\begin{grammarentry}{Abgrenzung}{Modalpartikel / Abtönungspartikel}
\GermanText{Als Satzadverb der Möglichkeit bedeutet vielleicht möglicherweise: Vielleicht regnet es morgen.}
\EnglishText{As a sentence adverb expressing possibility, vielleicht means perhaps: Perhaps it will rain tomorrow.}
\end{grammarentry}
\begin{exbox}\begin{itemize}
\ExampleItem{Vielleicht regnet es morgen.}{Perhaps it will rain tomorrow.}
\ExampleItem{Das war vielleicht eine Arbeit!}{What a lot of work that was!}
\end{itemize}\end{exbox}
\section{Lernhinweis \textnormal{(Learning tip)}}
\GermanText{Nicht die Wortform allein, sondern der Kontext entscheidet, ob vielleicht als Modalpartikel verwendet wird. Die englische Übersetzung kann die kommunikative Absicht statt einer wortwörtlichen Entsprechung ausdrücken.}
\EnglishText{The word form alone does not determine whether vielleicht is a modal particle. English translation may express its conversational intent instead of using a word-for-word equivalent.}
\section{Quellen \textnormal{(Sources)}}
\begin{itemize}
\item IDS Mannheim: \url{https://grammis.ids-mannheim.de/terminologie/2}
\item Duden: \url{https://www.duden.de/sprachwissen/sprachratgeber/Modalpartikeln}
\end{itemize}
\end{document}
